\section{C++ wrappers around the socket API}
\label{sec:wrappers}

Every backend in Section~\ref{sec:backends} is built from a small set of C++ classes
in \code{TcpSocket/}:
\begin{itemize}
  \item \code{TcpListener} and \code{TcpConnection}, which wrap one descriptor each;
  \item \code{PollSet} and \code{TcpConnectionPool}, built on \code{poll(2)};
  \item \code{EpollConnectionPool}, built on \code{epoll(7)};
  \item \code{TxPoller}, a write-only poller with a \code{poll} and an \code{epoll}
        back end.
\end{itemize}
This section describes how these classes call the C functions underneath, and which
rules of the C API shaped them. Table~\ref{tab:syscalls} lists every call and where
it is used.

\begin{table}[h]
\centering\small
\begin{tabularx}{\linewidth}{@{}>{\raggedright\arraybackslash}p{4.4cm}>{\raggedright\arraybackslash}X@{}}
\toprule
C function & Used in \\
\midrule
\code{socket}, \code{bind}, \code{listen}, \code{accept} & \code{TcpListener}; \code{SO\_REUSEADDR} set before \code{bind} \\
\code{socket}, \code{inet\_pton}, \code{connect} & \code{TcpConnection::connect}, always blocking \\
\code{setsockopt}, \code{fcntl} & \code{TcpConnection::apply\_options}: \code{SO\_SNDTIMEO}, \code{SO\_RCVTIMEO}, \code{TCP\_NODELAY}, \code{SO\_KEEPALIVE}, \code{O\_NONBLOCK} \\
\code{send}, \code{recv} & \code{TcpConnection::send}/\code{receive} (blocking, with retries) and \code{send\_once}/\code{receive\_once} (one call, caller's flags) \\
\code{shutdown}, \code{close} & \code{TcpConnection::shutdown} wakes other threads; only the destructor closes \\
\code{poll} & \code{PollSet::wait}, for \code{TcpConnectionPool} and the \code{poll} \code{TxPoller} \\
\code{epoll\_create1}, \code{epoll\_ctl}, \code{epoll\_wait} & \code{EpollConnectionPool} and the \code{epoll} \code{TxPoller} \\
\code{eventfd}, \code{read}, \code{write} & the wake-up descriptor of every poller \\
\bottomrule
\end{tabularx}
\caption{C functions and the wrappers that use them.}
\label{tab:syscalls}
\end{table}

\subsection{Owning a descriptor}

A descriptor is an \code{int} that must be closed exactly once, by exactly one
owner. \code{TcpConnection} and \code{TcpListener} make that an ownership rule of
the type:
\begin{itemize}
  \item the constructor takes the descriptor, and the destructor calls
        \code{close()} if it still holds one;
  \item the classes cannot be copied. A move takes the number with
        \code{std::exchange(other.fd\_, -1)}, so the moved-from object holds $-1$
        and its destructor does nothing;
  \item a connection shared between threads is held in a
        \code{std::shared\_ptr<TcpConnection>}. The last owner to let go closes the
        descriptor, so the number cannot be reused while another thread still uses
        it.
\end{itemize}

The C functions report errors by returning $-1$ and setting \code{errno}. The
wrappers translate \code{errno} at once into a small \code{SocketError} enumeration
(\code{EAGAIN} and \code{EWOULDBLOCK} become \code{WouldBlock}, \code{EINTR}
becomes \code{Interrupted}, and so on). They return it in a
\code{std::expected}, so the caller has to look at the error before it can use the
value. Only the \code{TcpListener} constructor throws, because a listener that
could not bind is not a usable object.

Addresses are built as a value-initialised \code{sockaddr\_in\{\}}, which also
zeroes the \code{sin\_zero} padding, and filled with \code{htons} and
\code{inet\_pton}. A failed \code{inet\_pton} is reported as a bad address before
any system call is made.

\subsection{Socket options}

\code{TcpConnection} stores its options and applies them in one place,
\code{apply\_options()}:
\begin{itemize}
  \item \code{SO\_SNDTIMEO} and \code{SO\_RCVTIMEO} take a \code{struct timeval}. A
        zero value means ``no timeout'', so the same code both sets and clears them.
  \item \code{TCP\_NODELAY} turns off Nagle's algorithm. Without it, a small frame
        could wait for the acknowledgement of the previous one. Every connection in
        the node sets it.
  \item \code{SO\_KEEPALIVE} lets the kernel notice a peer that has vanished without
        closing.
  \item \code{O\_NONBLOCK} is a flag on the open file, not a socket option. It is
        read with \code{fcntl(F\_GETFL)} and written back with \code{F\_SETFL}, so the
        other flags are kept.
\end{itemize}
The \code{with\_*()} setters only record the values, and the next \code{send} or
\code{receive} applies them. This is convenient for one thread. It is a data race
when a reader and a writer share the socket, because both would check and set the
``applied'' flag. Connections handed to a pool or a \code{TxPoller} are therefore
configured with \code{apply\_options\_now()} first; after that, the I/O calls only
read the flag.

\subsection{Reading and writing}

There are two pairs of I/O calls:
\begin{itemize}
  \item \code{send()} and \code{receive()} serve the blocking modes. They loop over
        \code{::send} and \code{::recv}, retry on \code{EINTR} and \code{EAGAIN} if
        retries were configured, and turn a \code{recv} that returns 0 into
        \code{ConnectionClosed}.
  \item \code{send\_once()} and \code{receive\_once()} are for the pollers. They make
        exactly one call with the flags the caller passes (\code{MSG\_DONTWAIT}),
        report \code{WouldBlock} as it is, and return a 0-byte read as 0, so the pool
        can treat it as end of stream.
\end{itemize}
Every \code{send} adds \code{MSG\_NOSIGNAL}. Without it, writing to a connection the
peer has closed raises \code{SIGPIPE}, which kills the whole process. With it, the
call simply fails with \code{EPIPE}.

A \code{send} may write only part of the buffer. The blocking outbound path wraps it
in \code{send\_all()}, which continues from where the previous call stopped. On a
blocking socket with \code{SO\_SNDTIMEO}, a send that waits too long fails with
\code{EAGAIN}, which the sender treats like a full buffer.

In \mode{threads} and \mode{pool-sender}, each reader thread blocks in
\code{receive()} with a 200\,ms \code{SO\_RCVTIMEO}. When the timeout expires,
\code{recv} fails with \code{EAGAIN}. The reader treats that as ``nothing yet'',
updates its heartbeat for the watchdog, and calls \code{recv} again.

\subsection{Ending a connection: \texttt{shutdown} before \texttt{close}}

Closing a descriptor that another thread is blocked on does not reliably wake that
thread. Worse, the number becomes free at once and can be given to a new connection
while the old thread still holds it. The wrappers therefore separate the two steps:
\begin{itemize}
  \item \code{TcpConnection::shutdown()} calls \code{shutdown(fd, SHUT\_RDWR)}. A
        thread blocked in \code{recv} on that socket then returns 0, as if the peer
        had closed. The descriptor stays valid, so no other connection can get its
        number. \code{ENOTCONN} is not treated as an error, because it only means
        the peer was already gone.
  \item \code{close()} is called only by the destructor, once no thread can use the
        socket any more.
\end{itemize}
The watchdog uses this to evict a blocking reader. The split modes use it so that a
failed write on the write thread shows up on the read thread as end of stream. The
pools call it whenever they drop a connection.

\code{accept()} on a listening socket cannot be woken this way. To stop a listener
blocked in \code{accept}, the node sets a stop flag and then connects to its own
listening port. \code{accept} returns, the loop sees the flag, and it closes the
wake-up connection.

\code{connect()} is always blocking (Section~\ref{sec:lessons}). The wrapper refuses
port~0 before calling it, because Linux does not reject a connection to port~0: the
SYN simply goes unanswered and the call blocks for about 130\,s of retries.

\subsection{The \texttt{poll} wrapper}
\label{sec:wrap-poll}

\code{poll()} reads an array of \code{struct pollfd} while it blocks, so no other
thread may change that array during the call. \code{PollSet} solves this with two
copies of the state:
\begin{itemize}
  \item a \emph{desired} table, \code{id}~$\to$~\code{pollfd}, which any thread may
        change under the pool's mutex (\code{upsert()} and \code{erase()});
  \item the \emph{kernel} array \code{fds\_}, which only the I/O thread touches. It
        is rebuilt from the desired table (\code{prepare()}), under the mutex, and
        only when that table has changed.
\end{itemize}
The I/O thread then releases the mutex and calls \code{::poll(fds\_, n, timeout)}.
The timeout is the time until the next sweep for idle connections. A call that fails
with \code{EINTR} is repeated.

Slot~0 of the array always holds an \code{eventfd} created with
\code{EFD\_NONBLOCK | EFD\_CLOEXEC}. When another thread adds a connection, queues a
write or closes the pool, it writes the 8-byte value 1 to that descriptor. The
blocked \code{poll} returns, and the I/O thread reads the counter back to zero and
rebuilds the array if needed.

On return, \code{for\_each\_ready()} walks the \code{revents} of every entry,
starting after the socket it served last, so a busy socket cannot starve the others:
\begin{itemize}
  \item \code{POLLOUT} with queued bytes: flush them with \code{send\_once};
  \item \code{POLLIN}: one \code{receive\_once} into the caller's buffer;
  \item \code{POLLERR}, \code{POLLHUP} or \code{POLLNVAL}: drop the connection and
        report it as closed.
\end{itemize}
\code{POLLOUT} is in the desired table only while a connection has queued bytes.
Otherwise every connection with free buffer space would be reported as writable on
every call.

\code{receive()} returns as soon as it has one chunk from one socket. The next call
makes a new \code{poll()}. With many sockets, every message read therefore costs one
scan of the whole array in the kernel and one walk over the \code{revents} in the
pool. This is the cost that grows with idle peers in S2.

\subsection{The \texttt{epoll} wrapper}

\code{EpollConnectionPool} uses the three \code{epoll} calls:
\begin{itemize}
  \item \code{epoll\_create1(EPOLL\_CLOEXEC)} creates the instance once, in the
        constructor. Its \code{eventfd} is registered first, with the token 0 in
        \code{epoll\_event.data.u64}.
  \item \code{epoll\_ctl} keeps the kernel's interest list up to date:
        \code{EPOLL\_CTL\_ADD} with \code{EPOLLIN} when a connection arrives,
        \code{EPOLL\_CTL\_MOD} to add or remove \code{EPOLLOUT} (only when the
        interest actually changes, which saves a system call), and
        \code{EPOLL\_CTL\_DEL} when it leaves. These calls are safe while the I/O
        thread is blocked in \code{epoll\_wait}, so no second copy of the state is
        needed.
  \item \code{epoll\_wait} fills a preallocated array of \code{max\_connections + 2}
        events.
\end{itemize}

Two details of the C API matter here:
\begin{itemize}
  \item \textbf{The event carries the connection id, not the fd.} Ids start at 1 and
        are never reused, so the token 0 can mean the \code{eventfd}, and an event
        for a connection removed in the meantime cannot reach a new connection with
        the same descriptor number.
  \item \textbf{\code{EPOLL\_CTL\_DEL} comes before \code{shutdown} and the release.}
        The kernel removes a descriptor from an epoll set by itself only when every
        reference to the underlying open file is closed. A connection shared with a
        \code{TxPoller} may outlive the pool's reference, so the pool removes it
        explicitly.
\end{itemize}

The pool is level-triggered: it does not set \code{EPOLLET}. A socket that still has
unread data is reported again by the next \code{epoll\_wait}, so the pool may read
one chunk and return. With edge triggering it would have to read until
\code{EAGAIN}, or it could miss data. Like the \code{poll} pool, \code{receive()}
serves one ready socket per call. Its fairness cursor indexes the events array of
the \emph{previous} call, whose contents are different each time, so it is not the
strict round robin of \code{PollSet}. This may contribute to \code{epoll}'s worse
tail at 1,024 busy peers (Section~\ref{sec:res-fanin}); we have not tested it.

\subsection{Writes through a pool, and the \texttt{TxPoller}}

All three pollers (\code{TcpConnectionPool}, \code{EpollConnectionPool} and
\code{TxPoller}) write in the same way:
\begin{enumerate}
  \item If nothing is queued for the connection, \code{send\_once(MSG\_DONTWAIT)}
        writes as much as the kernel accepts at once.
  \item Whatever the kernel did not accept is appended to the connection's
        \code{tx\_buffer}, and writability is armed: \code{POLLOUT} in the desired
        table, or \code{EPOLLOUT} through \code{epoll\_ctl}.
  \item When the socket becomes writable, the buffer is flushed until it is empty
        or the kernel says \code{EAGAIN} again. Once it is empty, writability is
        disarmed.
  \item A message that would push the buffer past 1\,MiB is refused with
        \code{WouldBlock}, and nothing of it is queued. This is the backpressure drop
        of Section~\ref{sec:arch}.
\end{enumerate}
The buffer keeps an offset of its first unsent byte, and it is compacted only once
that offset passes half the buffer, so appending stays cheap.

The \code{TxPoller} of the split modes runs the same algorithm on its own thread. It
never reads, and it watches a socket only while that socket has queued bytes, so its
set is usually empty. Unlike the read pools, it handles every ready descriptor
returned by one \code{poll} or \code{epoll\_wait} (up to 64 for \code{epoll}) before
it waits again.
